% !TEX program = xelatex
\documentclass[11pt]{ctexart}
\usepackage[a4paper,margin=2.4cm]{geometry}
\usepackage{booktabs}
\usepackage{longtable}
\usepackage{amsmath}
\usepackage{amssymb}
\usepackage{hyperref}
\hypersetup{colorlinks=true,linkcolor=blue!50!black,urlcolor=blue!50!black}
\usepackage{xcolor}
\usepackage{enumitem}
\setlist{nosep,leftmargin=1.6em}
\usepackage{caption}
\captionsetup{font=small,labelfont=bf}

\title{PyHost-Context：研究设计、探索结果、机制案例与效度限制\\[4pt]
\large 固定软件栈下 Python 测试的进程环境敏感性}
\author{ClawsGO Science Agent}
\date{2026-10-16}

\begin{document}
\maketitle
\vspace{-1.2em}
\begin{center}\small
本稿为\textbf{探索结果章节}的可写内容：方法、结果、机制案例与效度限制。所有数字由
\texttt{manuscript\_numbers.json} 从冻结证据重算，不手抄。独立验证组尚未执行，
全样本结论与摘要数值保持待定。
\end{center}

\tableofcontents
\bigskip

\section{研究设计与测量方法}

\subsection{研究问题与总体设计}

本研究考察一个受限的问题：在\textbf{代码、依赖、解释器与平台全部固定}的前提下，
只改变测试进程的启动环境，能否改变测试的可观察结果；如果能，变化经由哪条读取路径产生；
以及针对被观察到的变化，不同修复策略各自保留了什么。

设计上把三个问题对应到三个可分别作答的层次：\textbf{RQ1 变化}（在冻结条件下，
哪些（项目，条件）单位出现了可重复观察的结果差异）、\textbf{RQ2 机制}
（差异由哪一层读取环境、测试在哪一步建立了它的前置条件）、\textbf{RQ3 修复}
（策略消除被观察失败之后，还保留多少对原行为的覆盖）。

三条贯穿全篇的纪律：

\begin{enumerate}
\item \textbf{分母不可删}。没有观察到变化的单位、无法归因的单位、边界情况，一律留在分母里。
\item \textbf{三层证据分离}。“完整执行”“测试全部通过”“扰动按计划提交”是三个不同维度，
不合并成一个通过率。
\item \textbf{归因分层}。观察到的差异只有在满足约定条件后才升级到更高证据等级；
等级描述的是\textbf{支持什么陈述}，不是“缺陷确认”。
\end{enumerate}

\subsection{采样框与准入}

采样框取自公开的 PyPI 下载量排名快照（\texttt{top-pypi-packages} 30 天窗口），
按扫描深度 2304 逐条读取元数据，仅使用元数据，不执行任何候选代码。
在五个预设类别中各取至多 20 个合格项目作为抽取池，各类别合格数：
cli 38、data 108、dev\_test\_tools 157、text\_config 106、time\_i18n 20。

抽取后再经准入检查（可安装、可在冻结栈下跑通基线、sdist 字节可用哈希核验）。
最终冻结名单 26 项，按名次划分为\textbf{探索组 18 项}与
\textbf{留出验证组 8 项}（yaspin、about-time、leather、h5netcdf、flaky、text-unidecode、rich-rst、pytz-deprecation-shim）；
各类别实际入选数：cli 6、data 6、dev\_test\_tools 4、text\_config 6、time\_i18n 4；不足目标的部分如实记录（dev\_test\_tools 缺 2、time\_i18n 缺 2），
不替换、不补抽。验证组是\textbf{流行度偏移的留出}，不是随机独立同分布样本；
其扰动执行属于另行定义的阶段，本稿不含其结果。

\subsection{冻结软件栈}

解释器 CPython 3.11.9，pytest 7.2.0；
插件自动加载被关闭，入口点显式指定；缓存提供器禁用；
测试进程的外层超时为 240 秒。候选项目自身的测试在运行期不允许访问网络。
每轮的源码树在运行前按工件的哈希\textbf{逐字节恢复并复核}，随后才启动测试进程；
运行根下的 \texttt{home}/\texttt{tmp} 每轮清空重建并记录哈希。

\subsection{条件表与基线}

八个单键条件各自只改变一个环境键，其余键与基线 B0 完全一致（逐键断言，含仪器化键）：

\begin{table}[h]
\centering
\caption{冻结的八个单键条件与基线取值}
\small
\begin{tabular}{lllll}
\toprule
条件 & 环境键 & B0 & 扰动 & 说明 \\
\midrule
\texttt{K01\_TZ} & \texttt{TZ} & \texttt{UTC} & \texttt{Asia/Shanghai} & 时区 \\
\texttt{K02\_LC\_ALL} & \texttt{LC\_ALL} & 未设 & \texttt{C} & 区域设置类别 \\
\texttt{K03\_LANG} & \texttt{LANG} & \texttt{C.UTF-8} & \texttt{C} & 语言环境（可触发启动期 LC\_CTYPE 中介） \\
\texttt{K04\_HASH} & \texttt{PYTHONHASHSEED} & \texttt{0} & \texttt{42} & 哈希随机化种子 \\
\texttt{K05\_NO\_COLOR} & \texttt{NO\_COLOR} & 未设 & \texttt{1} & 彩色输出抑制 \\
\texttt{K06\_TERM} & \texttt{TERM} & \texttt{xterm-256color} & \texttt{dumb} & 终端类型 \\
\texttt{K07\_COLUMNS} & \texttt{COLUMNS} & 未设 & \texttt{40} & 输出宽度 \\
\texttt{K08\_LINES} & \texttt{LINES} & 未设 & \texttt{10} & 输出高度 \\
\bottomrule
\end{tabular}
\end{table}

基线 B0 不是“未设置任何变量”，而是\textbf{空白环境加显式白名单}：禁止未记录的宿主继承，
每个键以 未设／空串／设值 三态描述。这使“启动字典的差异恰好等于目标键”成为一个可断言的性质，
而不是一种期望。

\subsection{执行、顺序与测量控制}

每个（项目，条件）单位在\textbf{主矩阵}中重复 5 轮；另有\textbf{等执行预算的独立 B0 对照}，
每项目 45 轮。主矩阵 810 轮加对照 810 轮，
探索阶段共 1620 轮，另加恢复 12 轮与修复 28 轮，
合计 1660 次执行。可行性研究阶段的 336 次正式执行已另行封存，
\textbf{不并入本研究的分母}。

轮次顺序由调度指纹固定（对规范标识取哈希后按字典序执行），主矩阵轮与对照轮交错，
执行顺序本身写进交付件并由核验方重算。

每一轮留三层环境证据：\emph{启动字典}（交给子进程的全部环境键）、
\emph{初始化后快照}（解释器起来之后、收集之前）、
\emph{收集前快照}（测试进程内的运行期快照）。
三者的一致性构成“扰动按计划提交”的判据；允许的唯一例外是启动期已登记的
\texttt{LC\_CTYPE} 强制转换（仅在语言环境条件下适用，且需同时满足若干启动态前提）。

\subsection{记录规则与证据阶梯}

每轮为每个测试节点记录：原始结局、各阶段（setup／call／teardown）的原生状态、
节点身份指纹与失败签名。\textbf{测试身份变化是叠加标记，不是结果}：
若收集到的节点集合在轮次之间改变，该单位的\textbf{结果归因记为未解决}，
而不是把身份差异当作一次结果翻转。

证据阶梯 E0—E4 描述的是“支持什么陈述”：E4 的含义是
“在该项目、该键、该取值与该范围内支持变量归因”，
\textbf{不等于确认功能缺陷，也不等于安全消除}。

\subsection{恢复与修复端点}

对被观察到的翻转单位追加\textbf{恢复块}：块内固定顺序 参照$A_1$ → 扰动$B$ → 参照$A_2$，
两块相互独立，第二块的批次顺序为第一块的严格反序；判据要求
$A_1$ 与 $A_2$ 的全部记录一致、$B$ 相对 $A_1$ 的全部差异恰好等于签署的翻转集合。
两个合格块加块间调度门给出 E4。

\textbf{修复}只允许改测试、不允许改产品，且只在独立的派生工件上进行：冻结的行为变异施加在
派生源码树上，原始工件不被触碰。每个（单位 × 策略）判定三条并列的判据：
$A$ 被观察的失败在有效运行上下文中消失；$B$ 原有断言义务未被删除、跳过或放宽；
$C$ 同一枚变异探针仍在\textbf{精确的目标渲染断言}上于 \texttt{call} 阶段失败，
且失败差分里出现事前声明的偏差。另有共享参考轮与共享变异正控制轮，
两者不成立时该单位一律不得判为通过。

\section{探索结果}

\subsection{执行完整性}

探索阶段 1620 轮全部为\textbf{完整执行}：无超时、无启动字典越界、
无未登记的启动期改写（唯一登记例外是语言条件下的一次 \texttt{LC\_CTYPE} 转换，
且逐轮留证）。原始 JSONL 共 328,824,158 字节，畸形行 0 行。

返回码为 0 的 1250 轮，为 1 的 370 轮。
\textbf{返回码 1 是完整的会话并保留了测试失败观测，不是仪器故障}：
这些轮次中，项目自身的测试套件在基线环境里本来就有失败。
返回码 1 在对照轮 180 轮、主矩阵轮 190 轮，
按项目分布为 docutils（90）、flake8-bugbear（90）、python-dotenv（90）、rich-argparse（10）、x-wr-timezone（90）。

\subsection{单位级结果分层}

主分母是 144 个（项目，扰动条件）单位，即 18 个项目 ×
8 个单键扰动条件（B0 作为参照单独计入对照轮，不构成扰动单位）。
分层如下：

\begin{table}[h]
\centering
\caption{144 个扰动单位的结果分层（分母固定，未删任何一类）}
\begin{tabular}{lrl}
\toprule
层次 & 数量 & 说明 \\
\midrule
未观察到结果翻转 & 134 & 仅限冻结条件、取值与重复次数，\\
                  &              & 不能推成“该项目不存在环境假设” \\
测试身份未解决   & 8 & 全部来自同一项目（isodate），\\
                  &              & 无法把结果归因到同一节点，保留在分母 \\
达到 E4 的合格单位 & 2 & 同一项目的两个条件 \\
\midrule
单位总数（分母） & 144 & \\
\bottomrule
\end{tabular}
\end{table}

有翻转的单位共 10 个：其中 8 个因身份未解决
不能完成结果归因，2 个通过恢复复核达到 E4。
144 个单位的执行与启动态全部合规。

\subsection{两个达到 E4 的单位}

rich-argparse 在 \texttt{K05\_NO\_COLOR}（彩色输出抑制）下 22 个节点翻转，
在 \texttt{K07\_COLUMNS}（输出宽度）下 4 个节点翻转。
两组翻转在同一项目的 5 轮中稳定出现、与独立 B0 对照的噪声不相交、
且对照在全部翻转节点上都观测到同一结果。

两组节点重叠 4 个，合计 \textbf{22 个不同节点、1 个项目}。
恢复复核的 4 个块全部合格（$A_1/A_2$ 的完整身份、向量、阶段原生结局与
失败签名相等；$B$ 的全部差异恰为签署的翻转集合；有效起始源码哈希等于 E2 锚点）。

\textbf{计数纪律}：节点数不等于缺陷数，也不等于独立根因数，更不等于独立统计样本数。
22 个节点来自一个项目，不能读成 22 个独立发现。

\subsection{尚未解决的测试身份漂移}

另一个项目在\textbf{全部八个条件}下都出现收集到的测试身份发生变化：
某些参数化节点的标识在条件轮与对照轮之间不同，每轮收集总数不变。
该单位的逐测试结论因此无法在同一节点上比较，\textbf{结果归因记为未解决}。

必须区分的两种读法：这是\textbf{身份漂移}这一操作性观察；
本稿\textbf{不据此推断}参数表由库自身输出构造，也不把它算作八项独立发现或八项已确认的假阳性。
它保留在分母里，作为测量中必须报告的未解决项。

\subsection{修复端点}

修复阶段 28 个计划轮全部完成：每单位 14 轮 = 2 个共享控制 + 三策略各 4 格，
两单位合计 28 轮。每单位上限 15 轮、合计上限 30 轮，未用满的部分不补足。
2 个单位 × 3 条策略 = 6 个（单位×策略）组合中，
6 个满足全部判据（含共享参考复现、共享变异正控制与执行完整性）。

\textbf{该结果只在以下范围内成立}：本项目、本观察上下文、同一枚复合行为变异探针；
6 个通过组合不等于 6 个独立修复发现。

\section{机制与修复案例}

\subsection{机制一：依赖读取的颜色开关与未建立的测试前置条件}

观察到的差异是\textbf{表示层}的：同一断言下，条件轮的渲染输出缺少颜色码，正文文本不变。

读取路径可落到冻结源字节：
依赖 \texttt{rich} 在控制台构造时读取该环境键
（\texttt{console.py} 的 \texttt{Console.\_\_init\_\_}：\texttt{no\_color} 的取值源是
\texttt{\_environ.get("NO\_COLOR", "") != ""}，即\textbf{存在且非空}才抑制颜色，
空字符串不等价）；渲染时经去色处理（\texttt{\_render\_buffer} → \texttt{Segment.remove\_color}），
而去色只清颜色、保留粗体等属性。

产品侧的耦合方式是\textbf{惰性}的：formatter 的 \texttt{console} 属性在首次访问时才构造
\texttt{r.Console()}，因此拿到的是\textbf{访问当时}的进程环境。

测试侧的问题不在断言本身，而在前置条件：
需要颜色的测试用一个 fixture 把强制着色变量设为 1，而该 fixture 只做“增加”，
\textbf{从不建立被抑制键的缺席}。当被抑制键在启动环境里存在时，产品行为是合规的，
而测试期望的字面输出仍然是有颜色的——两者冲突，冲突出在测试的环境前提上。

两轴分类（产品契约符合性／测试是否显式建立其断言的环境）在该单位上取值为
\texttt{conformant}／\texttt{ambient}，差异标签为表示层变化。
契约符合性\textbf{只限定于观察到的这条颜色行为}，不是整产品认证。

\subsection{机制二：标准库在收集期捕获宽度，会话 fixture 设置过晚}

第二个单位的差异同样只体现在排版上：文本内容不变，折行位置不同。

读取宽度的是\textbf{标准库}：\texttt{optparse} 的 \texttt{HelpFormatter.\_\_init\_\_}
在宽度为 \texttt{None} 时读取环境键，缺失取 80，随后减 2。
产品按该宽度属性折行。

测试侧的时机问题是可静态定位的：被观察的测试在\textbf{模块导入时的参数表}里构造
formatter 对象，也就是在\textbf{收集期}；而把环境宽度设为 100 的 fixture 是会话作用域，
设置在\textbf{收集之后}生效。因此宽度在 fixture 生效之前就已经被固定下来了。
被测测试所依据的字面布局对应一个具体宽度，而它从未把这个宽度建立起来。

本例说明一个与上一例不同的作用路径：前者是依赖在运行期读环境，后者是标准库在\textbf{收集期}
读环境——两者都属于“测试对环境建立不足”，但读取层、时机与修复方式都不同。

\subsection{三条修复策略与各自观察到的覆盖范围}

修复阶段比较三种策略：\textbf{全局固定}（在会话开始前把相关环境键固定下来）、
\textbf{局部隔离}（只在需要该前置条件的测试内建立条件）、
\textbf{显式参数化}（把环境取值写成测试参数，按各上下文取期望值）。
三策略在两单位上都消除了被观察的失败，并在精确目标渲染断言上保留了对同一枚
cyan$\rightarrow$magenta 变异探针的敏感性：

\begin{table}[h]
\centering
\caption{三策略 × 两单位的端点结果（目标命中数为该格命中的节点数）}
\small
\begin{tabular}{lllrrl}
\toprule
单位 & 策略 & 有效运行上下文 & 有效上下文命中 & B0 命中 & 状态 \\
\midrule
\texttt{NO\_COLOR} & R1 & Z & 1 & 1 & passed \\
\texttt{NO\_COLOR} & R2 & B & 1 & 1 & passed \\
\texttt{NO\_COLOR} & R3 & B & 1 & 1 & passed \\
\texttt{COLUMNS} & R1 & Z & 4 & 4 & passed \\
\texttt{COLUMNS} & R2 & B & 4 & 4 & passed \\
\texttt{COLUMNS} & R3 & B & 8 & 8 & passed \\
\bottomrule
\end{tabular}
\end{table}

\textbf{本数据不能给出策略优越性排序}：三者在两单位上都满足约定判据。
有信息量的差异在\textbf{覆盖范围}，而不在“是否通过”：

\begin{itemize}
\item \textbf{全局固定}把有效运行上下文变成 Z（该键被固定为一个取值）。
它让问题不再出现，但产品原本支持的另一种行为不再有任何断言覆盖；
报告不得声称“修复后的套件仍在原扰动环境中通过”。
\item \textbf{局部隔离}只在需要颜色的测试内建立前置条件；
被抑制键在场时产品的行为仍无断言。
\item \textbf{显式参数化}把环境取值纳入测试，两个取值都跑到并按各自上下文取期望值，
覆盖范围最广；但参数化套件里真正携带颜色、因而\textbf{能}检出颜色类变异的分支只有一个，
另一个分支只是不再因该键而失败。不得由整个参数化套件失败推导每个分支都能检测该探针。
\end{itemize}

宽度的第三种策略同样需要收窄表述：它验证的是两种环境取值下的\textbf{宽度读取}
（显式宽度的渲染，加上新建 formatter 的宽度属性），\textbf{没有}验证两个宽度下的完整自动排版。

\section{效度限制}

\subsection{选择与配额}

采样受三重限制，逐条都会影响可推广性。其一，抽取池按下载量名次与五个预设类别配额构成，
不是从 Python 生态的随机样本；其二，准入需要候选项目在冻结栈下可安装、可跑通基线，
这本身就是一层筛选，排除了装不上、跑不动的项目；
其三，名次前段的探索组与后段的留出验证组之间存在流行度偏移，
验证组是\textbf{留出}而不是随机抽样。因此本项目出现的比例不能读成生态发生率。

准入阶段另有若干类别未能凑满目标数（dev\_test\_tools 缺 2、time\_i18n 缺 2），如实记录、不替换、不补抽。

\subsection{条件与取值}

八个条件各自只取\textbf{一个}扰动取值，且只改一个键。
因此本设计能回答“该键在该取值上是否改变了可观察结果”，
不能回答“该键在整个取值范围内的影响曲线”。
未设与空串在多数条件下被区分对待，但并非全部条件都取了这一对取值。
五个类别的配额也意味着某些类别（如时序与本地化）在样本中占比很小。

\subsection{重复次数与筛查功效}

每单位 5 轮只用于\textbf{筛查}，不作为统计功效依据。
在单次失败概率为 0.10 的假设下，五次至少出现一次失败的概率约为 0.40951；
也就是说，五轮设计对低概率不稳定的检出能力有限，
“未观察到翻转”不能被读成“不存在环境假设”。

\subsection{独立对照与噪声底}

独立 B0 对照每项目 45 轮，用于估计同一配置重复执行时的噪声底。
本批只有一个项目在对照中出现了节点级噪声，而它正是出现身份漂移的那个项目；
翻转与噪声相交的单位一律不计为扰动证据。
但对照本身也只覆盖了有限轮次，不能排除更低频率的噪声。

\subsection{未解决项与缺失证据}

三处必须照录的局限。

第一，\textbf{测试身份漂移未解决}：该单位的节点集合在轮次之间改变，
无法把结果归因到同一节点；本稿只把它记为未解决，不推断其成因，
也不把它算作发现。

第二，\textbf{三份诊断快照缺失}：绘制最小复现时使用的三份渲染快照在交付前被采集脚本
连同临时树一并删除。同一轮运行的逐节点比较记录仍在，
但\textbf{不能声称由缺失的快照重建了首次或完整输出的等价性}。不回填、不重跑。

第三，\textbf{依赖闭包未重新认证}：只交付了少量关键依赖源文件用于读取链核对，
不等于重新构建或重新认证了整个冻结环境。

此外，修复端点所依据的行为变异只有\textbf{一枚}复合样式变异。
杀死单枚变异只支持“对该探针仍敏感”，不能证明全部断言义务、全部支持行为或全部环境取值
都得到了保留；差分被截断、看不出颜色变化的失败不作为偏差证据。

\subsection{归因边界}

E4 的含义被严格限定为“在该项目、该键、该取值与该范围内支持变量归因”。
本稿\textbf{不确认功能缺陷}，不主张安全消除，也不主张该现象在生态中普遍存在。
两个 E4 单位属于\textbf{同一个项目}，重叠 4 个节点、合计 22 个不同节点，
机制描述为两条但\textbf{不主张统计独立性}。
把机制数、节点数与项目数三者混同，是本设计最容易出现的误读。

\subsection{尚待完成的工作}

冻结样本共 26 项，本研究目前报告探索组 18 项。
留出验证组 8 项（yaspin、about-time、leather、h5netcdf、flaky、text-unidecode、rich-rst、pytz-deprecation-shim）尚未执行，其扰动执行属于另行定义的阶段，
本轮未获授权；其主矩阵与独立 B0 对照轮次按同一设计待定。
原计划的工具原生策略对照、全量成本统计与端点汇总亦未由现有证据支持，不写成已完成结果。

因此，\textbf{验证结束之前，全样本结论与摘要数值应保持待定}；
本稿报告的探索结果可以先行成为完整章节，但不承担全样本层面的结论。

\bigskip
\noindent\rule{\linewidth}{0.4pt}

\smallskip
\noindent\small
\textbf{数字来源}：本节全部计数由 \texttt{extension/tooling/build\_manuscript\_numbers.py}
从冻结证据重算并写入 \texttt{extension/analysis/manuscript\_numbers.json}，
再由 \texttt{build\_manuscript\_sections.py} 排版；正文不含手写数字。
证据等级、修复判据与边界表述沿用已签署的评审结论。

\end{document}
